{\begingroup



\section{Symplectic Aspects of Toric Varieties}

Let \(X\) be a smooth projective toric variety. Since \(X\) is projective, its
underlying complex manifold admits a Kähler form. In particular, \(X\) may be
viewed as a symplectic manifold. The purpose of this section is to recall the
basic symplectic notions needed for this perspective and to explain, at the level
of definitions, how the symplectic structure on \(X\) is obtained from the
algebraic data of an ample toric line bundle.

We shall also fix notation and introduce several related concepts that will be
used later in the discussion of toric coherent-constructible correspondence. To
keep the exposition concise, we will state the standard facts used in this section
without proof.

In this chapter, we will use both the Zariski topology and the usual analytic
topology when needed, and we will regard \(X\) both as an algebraic variety and
as its associated complex manifold; when no confusion can arise, we will not
explicitly distinguish between these viewpoints or topologies.






\subsection{Symplectic Manifolds}

Let \(X\) be a smooth real manifold. A \textbi{symplectic form} on \(X\) is a
closed non-degenerate two-form
\[
\omega\in \Omega^2(X),
\]
that is,
\[
d\omega=0,
\]
and for every point \(x\in X\), the bilinear form
\[
\omega_x:T_xX\times T_xX\to \mathbb R
\]
is non-degenerate. Equivalently, the map
\[
T_xX\longrightarrow T_x^*X,\qquad
v\longmapsto \omega_x(v,-)
\]
is an isomorphism for every \(x\in X\). A pair \((X,\omega)\) is called a
\textbi{symplectic manifold}.

A diffeomorphism
\[
\phi:X\to X
\]
is called a \textbi{symplectomorphism} if it preserves the symplectic form:
\[
\phi^*\omega=\omega.
\]
Thus a symplectomorphism is an automorphism of the symplectic manifold
\((X,\omega)\).

Let \(G\) be a Lie group acting smoothly on \(X\). We say that the action is a
\textbi{symplectic action} if every element \(g\in G\) acts by a symplectomorphism,
that is,
\[
g^*\omega=\omega
\]
for all \(g\in G\).
We denote the action of $G$ by
\[
\alpha:G\times X\to X,\qquad (g,x)\mapsto g\cdot x.
\]
For each point \(x\in X\), the action determines an orbit map
\[
\alpha_x:G\to X,\qquad g\mapsto g\cdot x.
\]
Let $\mathfrak{g}$ be the Lie algebra of \(G\). For \(\xi\in\mathfrak g\), the corresponding
one-parameter subgroup is
\[
t\longmapsto \exp(t\xi).
\]
Differentiating the orbit map along this one-parameter subgroup gives a tangent
vector at \(x\):
\[
\xi_X(x)
=
\left.\frac{d}{dt}\right|_{t=0}
\exp(t\xi)\cdot x
\in T_xX.
\]
As \(x\) varies, this defines a vector field
\[
\xi_X\in \Gamma(TX).
\]
The assignment
\[
\mathfrak g\to \Gamma(TX),\qquad
\xi\mapsto \xi_X
\]
is called the \textbi{infinitesimal action} of \(G\) on \(X\). The vector field
\(\xi_X\) is called the \textbi{fundamental vector field} generated by \(\xi\).

A symplectic \(G\)-action is called \textbi{Hamiltonian} if there exists a smooth
map
\[
\mu:X\to\mathfrak g^*
\]
such that, for every \(\xi\in\mathfrak g\),
\[
d\langle \mu,\xi\rangle
=
\iota_{\xi_X}\omega .
\]
Here
\[
\langle \mu,\xi\rangle:X\to\mathbb R
\]
denotes the function obtained by pairing \(\mu(x)\in\mathfrak g^*\) with
\(\xi\in\mathfrak g\), and \(\iota_{\xi_X}\omega\) denotes contraction of
\(\omega\) with the vector field \(\xi_X\).
The map
\[
\mu:X\to\mathfrak g^*
\]
is called a \textbi{moment map}. The triple \((X,\omega,\mu)\) is called a
\textbi{Hamiltonian \(G\)-space}.

The moment map is not unique in general. If \(\mu:X\to\mathfrak g^*\) is a moment
map, then \(\mu+c\) is also a moment map for any constant element
\(c\in\mathfrak g^*\). Thus a moment map is determined only up to addition of a
constant element of \(\mathfrak g^*\).

A moment map is called \textbi{\(G\)-equivariant} if
\[
\mu(g\cdot x)=\operatorname{Ad}^*_g \mu(x)
\]
for all \(g\in G\) and \(x\in X\), where \(\operatorname{Ad}^*\) is the coadjoint action, namely the dual action to
the adjoint representation
\[
\operatorname{Ad}:G\to \operatorname{GL}(\mathfrak g).
\]


\subsection{Kähler Manifolds and Complex Structure}
Let \(X\) be a smooth manifold, and let \(D\subset TX\) be a smooth subbundle of
the tangent bundle. We call \(D\) a \textbi{distribution} on \(X\). 

The distribution \(D\) is called \textbi{integrable} if for every point \(x\in X\),
there exists an immersed submanifold \(Y\subset X\) passing through \(x\) such
that
\[
T_yY=D_y
\]
for every \(y\in Y\). Such a submanifold \(Y\) is called an \textbi{integral manifold} of
\(D\).

An \textbi{almost complex structure} on \(X\) is a
bundle endomorphism
\[
J:TX\to TX
\]
such that
\[
J^2=-\operatorname{id}_{TX}.
\]

If \(J\) is integrable, then \(X\) is a complex manifold, and \(J\) is called its
\textbi{complex structure}.

Suppose now that \(X\) is a complex manifold with complex structure \(J\). 

A Riemannian metric \(g\) on \(X\) is called \textbi{Hermitian} if it is compatible with
\(J\), namely
\[
g(Jv,Jw)=g(v,w)
\]
for all tangent vectors \(v,w\). Associated to such a Hermitian metric is a real
two-form
\[
\omega(v,w)=g(Jv,w).
\]
This two-form is called the \textbi{fundamental form} of the Hermitian metric.

The Hermitian metric \(g\) is called  \textbi{Kähler} if
\[
d\omega=0.
\]
In this case \(\omega\) is called a \textbi{Kähler form}, and the triple
\[
(X,J,\omega)
\]
is called a \textbi{Kähler manifold}.

A Kähler form is in particular a symplectic form. Indeed, it is a closed
non-degenerate real two-form. Thus every Kähler manifold is naturally a
symplectic manifold. Conversely, on a complex manifold, a symplectic form
compatible with the complex structure determines a Kähler metric by
\[
g(v,w)=\omega(v,Jw).
\]
\subsection{Hermitian Line Bundles and Curvature}
We recall a standard construction from complex differential geometry. Let \(X\)
be a complex manifold, and let $L$
be a holomorphic line bundle. A \textbi{Hermitian metric} on \(L\) is a smoothly varying
Hermitian inner product
\[
h_x:L_x\times L_x\to \mathbb C
\]
on the fibers of \(L\). If \(s\) is a local non-vanishing holomorphic frame of
\(L\), then the metric is determined locally by the positive smooth function
\[
h(s,s)>0.
\]

There is a unique connection \(\nabla_h\) on \(L\), called the \textbi{Chern connection},
which is compatible with both the holomorphic structure and the Hermitian metric.
The compatibility with the holomorphic structure means that
\[
\nabla_h^{0,1}=\bar\partial_L,
\]
while compatibility with the metric means that
\[
d\,h(s_1,s_2)
=
h(\nabla_h s_1,s_2)+h(s_1,\nabla_h s_2)
\]
for smooth local sections \(s_1,s_2\) of \(L\).
The \textbi{curvature} of the Chern connection is
\[
F_h=\nabla_h^2.
\]
Since \(L\) is a line bundle, \(F_h\) may be regarded as a complex-valued
two-form on \(X\).
Locally, if \(s\) is a holomorphic frame and
\[
h(s,s)=e^{-\varphi},
\]
then the curvature is expressed in terms of the local potential \(\varphi\). With
the usual convention, the associated real \((1,1)\)-form is
\[
\omega_h=\sqrt{-1}\,\partial\bar\partial \varphi.
\]
Equivalently, up to the chosen sign convention, one may write
\[
\omega_h=\sqrt{-1}F_h.
\]

If \(L\) is positive, then \(\omega_h\) is a Kähler form. Thus a positive
Hermitian metric on a positive holomorphic line bundle produces a symplectic
form on the underlying smooth manifold.



\subsection{Algebraic Tori and Compact Real Tori}

Let \(X\) be a smooth projective toric variety of complex dimension
\(n\). 
The algebraic torus \(T\) is a complex algebraic group. It contains a distinguished
compact real Lie subgroup, namely its maximal compact real subtorus
\[
T_{\mathbb R}\cong (S^1)^n.
\]
Concretely, under the identification \(T\cong(\mathbb C^*)^n\), this subgroup is
\[
T_{\mathbb R}
=
\{(z_1,\dots,z_n)\in(\mathbb C^*)^n
\mid |z_1|=\cdots=|z_n|=1\}.
\]
The \textbi{compact torus action} 
\[
T_{\mathbb R}\times X\longrightarrow X.
\]
is obtained by restricting the algebraic \(T\)-action on \(X\) to the subgroup
\(T_{\mathbb R}\subset T\).

The real vector space \(N_{\mathbb R}\) is naturally identified with the Lie
algebra of the compact torus \(T_{\mathbb R}\):
\[
N_{\mathbb R}\cong \operatorname{Lie}(T_{\mathbb R}).
\]
Dually, we have
\[
M_{\mathbb R}\cong \operatorname{Lie}(T_{\mathbb R})^*.
\]
Under these identifications, the \textbi{exponential map of the compact torus} is
\[
\exp:T_{\mathbb R}^{\mathrm{Lie}}\cong N_{\mathbb R}\longrightarrow T_{\mathbb R},
\]
and its kernel is the lattice \(N\). Hence
\[
T_{\mathbb R}\cong N_{\mathbb R}/N.
\]
Similarly, the \textbi{dual compact torus} is
\[
T_{\mathbb R}^{\vee}\cong M_{\mathbb R}/M.
\]






\subsection{Kähler Structures from Ample Toric Line Bundles}

Let \(X\) be a smooth projective toric variety over \(\mathbb C\), and let
\[
L=\mathcal O_X(D)
\]
be an ample \(T\)-linearized line bundle on \(X\). Passing to the analytification,
we regard \(L\) as a holomorphic line bundle over the complex manifold
\(X\).
Since \(L\) is ample, one can choose a positive smooth Hermitian metric \(h\) on
\(L\). Let \(\nabla_h\) be the associated Chern connection and let
\[
F_h=\nabla_h^2
\]
be its curvature. With a fixed sign convention, the form
\[
\omega_D=\sqrt{-1}F_h
\]
is a Kähler form on \(X\). In particular, \(\omega_D\) is a closed
non-degenerate real two-form, so
\[
(X,\omega_D)
\]
is a symplectic manifold.

Moreover, since the compact real torus $T_{\mathbb R}$
is compact, we may average the Hermitian metric \(h\) over \(T_{\mathbb R}\).
Thus \(h\), and hence \(\omega_D\), may be chosen to be \(T_{\mathbb R}\)-invariant.
Therefore the algebraic data $(X,\mathcal O_X(D))$
gives rise to a \(T_{\mathbb R}\)-invariant Kähler, hence symplectic, manifold $(X,\omega_D).$

Since \(\omega_D\) may be chosen to be \(T_{\mathbb R}\)-invariant, the compact
torus action on \(X_\Sigma\) is symplectic. Using the identifications
\[
\operatorname{Lie}(T_{\mathbb R})\cong N_{\mathbb R},
\qquad
\operatorname{Lie}(T_{\mathbb R})^*\cong M_{\mathbb R},
\]
a moment map for this action is a smooth map
\[
\mu_D:X_\Sigma\to M_{\mathbb R}
\]
such that
\[
d\langle \mu_D,\xi\rangle=\iota_{\xi_X}\omega_D
\]
for every \(\xi\in N_{\mathbb R}\). Thus the pair
\((X_\Sigma,\mathcal O_X(D))\) produces a Hamiltonian
\(T_{\mathbb R}\)-space
\[
(X_\Sigma,\omega_D,\mu_D).
\]
The image of \(\mu_D\) will be discussed in the next section.


\subsection{Moment Polytope}

Let \(X\) be a smooth projective toric variety. 
Let
\[
D=\sum_{\rho\in\Sigma(1)}c_\rho D_\rho
\]
be a \(T\)-invariant Cartier divisor, where \(D_\rho\) is the torus-invariant
divisor corresponding to the ray
\[
\rho=\mathbb R_{\geq 0}v_\rho
\]
and \(v_\rho\in N\) is the primitive generator of \(\rho\). Assume that
\[
L=\mathcal O_X(D)
\]
is ample.

\begin{definition}[Moment polytope]
The polytope associated to \(D\) is
\[
\Delta_D
=
\left\{
m\in M_{\mathbb R}
\ \middle|\
\langle m,v_\rho\rangle\geq -c_\rho
\text{ for all }\rho\in\Sigma(1)
\right\}.
\]
\end{definition}
The main fact we need is the following.
\begin{theorem}
With the standard normalization determined by the divisor \(D\), the image of the
moment map
\[
\mu_D:X_\Sigma\to M_{\mathbb R}
\]
associated to the ample line bundle \(\mathcal O_X(D)\) is precisely the
polytope \(\Delta_D\):
\[
\mu_D(X_\Sigma)=\Delta_D.
\]
\end{theorem}

\begin{proof}
We use the projective embedding determined by the complete linear system of
\(L=\mathcal O_X(D)\). By the standard description of global sections of toric
line bundles, we have
\[
H^0(X,\mathcal O_X(D))
=
\bigoplus_{m\in \Delta_D\cap M}\mathbb C\cdot \chi^m .
\]
Thus the lattice points of \(\Delta_D\) give a basis of global sections. Writing
\[
\Delta_D\cap M=\{m_0,\dots,m_s\},
\]
the complete linear system defines a \(T\)-equivariant projective embedding
\[
\Phi_D:X_\Sigma\longrightarrow \mathbb P^s.
\]
On the dense torus \(T\subset X_\Sigma\), this map is given by
\[
t\longmapsto
[\chi^{m_0}(t):\cdots:\chi^{m_s}(t)].
\]

We endow \(\mathbb P^s\) with the Fubini--Study Kähler form and pull it back to
\(X_\Sigma\). The corresponding moment map, restricted to the dense torus, is
\[
\mu_D(t)
=
\frac{
\sum_{m\in\Delta_D\cap M} m\,|\chi^m(t)|^2
}{
\sum_{m\in\Delta_D\cap M} |\chi^m(t)|^2
}.
\]
Therefore \(\mu_D(t)\) is a weighted average of the lattice points
\(m\in \Delta_D\cap M\), with positive real weights
\[
|\chi^m(t)|^2.
\]
Hence
\[
\mu_D(T)\subset \operatorname{Conv}(\Delta_D\cap M)=\Delta_D.
\]
Since \(T\) is dense in \(X_\Sigma\) and \(\mu_D\) is continuous, it follows that
\[
\mu_D(X_\Sigma)\subset \Delta_D.
\]

It remains to see that every point of \(\Delta_D\) occurs in the image. On the
dense torus, write
\[
t=\exp(y+i\theta),
\qquad y\in N_{\mathbb R},\quad \theta\in N_{\mathbb R}/N.
\]
Then
\[
|\chi^m(t)|^2=e^{2\langle m,y\rangle}.
\]
Thus the moment map on the open orbit can be written as
\[
\mu_D(y)
=
\frac{
\sum_{m\in\Delta_D\cap M}m\,e^{2\langle m,y\rangle}
}{
\sum_{m\in\Delta_D\cap M}e^{2\langle m,y\rangle}
}.
\]
Equivalently,
\[
\mu_D(y)
=
\frac{1}{2}\nabla_y
\log\left(
\sum_{m\in\Delta_D\cap M}e^{2\langle m,y\rangle}
\right).
\]
The gradient of the strictly convex function
\[
y\longmapsto
\log\left(
\sum_{m\in\Delta_D\cap M}e^{2\langle m,y\rangle}
\right)
\]
maps \(N_{\mathbb R}\) onto the interior of the convex hull of the weights
\(\Delta_D\cap M\), namely onto \(\Delta_D^\circ\). Hence
\[
\Delta_D^\circ\subset \mu_D(T).
\]
Taking closures gives
\[
\Delta_D\subset \mu_D(X_\Sigma).
\]
Combining the two inclusions, we obtain
\[
\mu_D(X_\Sigma)=\Delta_D.
\]
\end{proof}

Thus the algebraic datum of the ample toric line bundle \(\mathcal O_X(D)\)
determines not only a Kähler form \(\omega_D\), but also a canonical convex
polytope \(\Delta_D\), namely the image of the associated moment map.









\endgroup}